\begin{theorem}[Unitary equivalence of closure states on two width-one stripes]%
    \label{thm:peps_torus_stripe_local_equivalence}
    \lean{TNLean.PEPS.mem_compl_torusStripeInterior_iff_coordinate_zero,
        TNLean.PEPS.IsGIsometric.torusClosureSuperpositionCut_stripe_local_equivalence}
    \leanok
    \uses{thm:peps_torus_rectangle_realization,
        lem:peps_torus_rectangle_cut_connectivity,
        thm:peps_torus_simply_connected_physical_local_equivalence,
        def:peps_normalized_pure_cut_matrix}
    Consider an original regular $G$-isometric tensor on a native
    torus whose two periods are at least three. Let $R$ be the
    nonwrapping rectangle of sites with both coordinates nonzero.
    Its complement is the union of two width-one stripes wrapping
    around the torus:
    \begin{align}
        S & =\{(x,y):x=0\lor y=0\}.
        \label{eq:peps_width_one_stripes}
    \end{align}
    For any two nonzero coherent superpositions of commuting torus
    closures, let $A$ and $B$ be their physical coefficient matrices
    across $R\mid S$. There exists a unitary $U$ on the physical
    factor of $S$ satisfying
    \begin{align}
        \widehat B & =\widehat A U^{\mathsf T}.
        \label{eq:peps_stripe_physical_unitary}
    \end{align}
    Thus $U$ transforms the first normalized physical state into the
    second while acting only on the stripes
    (\ref{eq:peps_width_one_stripes}). This is the closure-vector
    specialization of \cite[Theorem~6.7]{Schuch2010PEPS} to regular
    native tori with periods at least three. Identification of the
    closure span with a parent-Hamiltonian ground space is separate.
\end{theorem}

\begin{proof}\leanok
    The coordinate-value bounds identify $R$ with the rectangle
    starting at $(1,1)$ with side lengths one less than the two periods.
    Its induced graph is connected, and
    \ref{thm:peps_torus_rectangle_realization} gives simple
    connectedness of its actual closed-cell realization.
    Apply \ref{thm:peps_torus_simply_connected_physical_local_equivalence}
    to obtain (\ref{eq:peps_stripe_physical_unitary}).
\end{proof}
